% ===== 1 引言与问题定义 =====
\section{引言与问题定义}
\label{sec:introduction}

\subsection{背景与竞赛任务}
\label{sec:task-background}

骨料是混凝土中体积分数最高的组成材料之一，其粒径级配与颗粒形貌直接影响混凝土拌合物的堆积结构、工作性能以及硬化后的力学性能。因此，对骨料粒径分布、颗粒形状及异常物进行稳定检测，是混凝土生产质量控制中的基础环节。工程中通常依据 \texttt{GB/T 14685}、\texttt{JGJ 52} 以及 \texttt{ASTM C136} 等标准，通过逐级机械筛分和称重获得不同筛孔区间内的质量占比，并进一步计算累计通过率及 $D_{10}$、$D_{50}$、$D_{90}$ 等统计量\cite{gb14685-2022,jgj52-2006}。该方法具有明确的物理意义和较好的可重复性，但需要取样、筛分和称重等离线操作，检测周期与人工参与程度较高，难以直接形成面向生产过程的快速反馈。

本赛题将上述工程需求转化为自然堆积或传送带场景下的非人工颗粒筛分任务，要求参赛系统完成颗粒精细识别、粒径分析、形状分类和异常检测，并在有限的现场采集与调试时间内给出可用于答辩和评价的筛分结果。与一般目标检测任务不同，比赛最终以标准筛分结果作为粒径评价依据，因此系统不仅需要回答“图像中有哪些颗粒”，还需要回答“这些可见颗粒对应怎样的工程筛分级配”。表\ref{tab:competition-tasks}概括了赛题要求及本文对应的系统输出。

\begin{table}[htbp]
\centering
\caption{赛题要求与本文系统输出的对应关系}
\label{tab:competition-tasks}
\small
\begin{tabular}{>{\raggedright\arraybackslash}p{2.2cm}>{\raggedright\arraybackslash}p{5.0cm}>{\raggedright\arraybackslash}p{6.0cm}}
\toprule
赛题要求 & 需要解决的问题 & 本文主要输出 \\
\midrule
精细识别 & 区分颗粒与背景，并分离相互接触的颗粒 & 颗粒实例掩膜 $\{M_i\}_{i=1}^{N}$ \\
粒径分析 & 将图像尺寸恢复为物理尺寸，并统计整体级配 & 长短轴、筛分等效粒径、$D_{10}/D_{50}/D_{90}$、各筛级质量占比 \\
形状分类 & 描述针片、圆形、棱角等颗粒投影形貌 & 长宽比、圆度、凸度及形貌类别 \\
异常检测 & 识别明显超限颗粒及疑似非骨料物质 & 超限、疑似泥团及其他异常标签 \\
现场约束 & 在规定时间内完成采集、调试、推理与结果输出 & 自动化处理流程、标准化图表与场景级报告 \\
\bottomrule
\end{tabular}
\end{table}

由此可见，本赛题的目标并不是简单以视觉模型替代人工观察，而是以图像为输入建立一条能够逼近标准筛分语义的自动化测量链。要实现这一目标，首先需要识别自然堆积图像与机械筛分之间存在的关键差异。

\subsection{自然堆积视觉筛分的核心挑战}
\label{sec:problem-challenges}

从一幅自然堆积骨料图像得到标准筛分意义下的质量级配，需要连续完成“图像像素$\rightarrow$独立颗粒$\rightarrow$物理几何$\rightarrow$筛分粒径$\rightarrow$质量分布”的跨层映射。该过程至少包含三个相互耦合的困难：实例边界不确定、像素尺度缺乏直接物理意义，以及二维视觉统计与机械筛分统计口径不一致。任一环节产生的偏差都会沿处理链传播，并最终表现为粒级质量占比或特征粒径的误差。图\ref{fig:three-gaps}给出了本文对该问题的分层理解。

\begin{figure}[htbp]
\centering
\begin{tikzpicture}[
  node distance=1.0cm,
  block/.style={rectangle, draw, rounded corners, minimum width=2.75cm, minimum height=0.95cm, align=center, fill=blue!6, font=\small},
  gap/.style={rectangle, draw, dashed, rounded corners, minimum width=2.45cm, minimum height=0.72cm, align=center, fill=orange!8, font=\footnotesize},
  arrow/.style={->, thick}
]
\node[block] (img) {自然堆积图像\\$I$};
\node[gap, right=0.75cm of img] (g1) {实例感知鸿沟\\粘连、遮挡、阴影};
\node[block, right=0.75cm of g1] (inst) {独立颗粒\\$\{M_i\}$};
\node[gap, below=0.95cm of inst] (g2) {物理尺度鸿沟\\pixel $\rightarrow$ mm};
\node[block, left=0.75cm of g2] (geo) {物理几何\\$a_i,b_i,A_i,\ldots$};
\node[gap, left=0.75cm of geo] (g3) {筛分语义鸿沟\\2D 数量 $\rightarrow$ 质量级配};
\node[block, below=0.95cm of g3] (grade) {标准筛分输出\\$D_p$ 与粒级质量占比};
\draw[arrow] (img) -- (g1);
\draw[arrow] (g1) -- (inst);
\draw[arrow] (inst) -- (g2);
\draw[arrow] (g2) -- (geo);
\draw[arrow] (geo) -- (g3);
\draw[arrow] (g3) -- (grade);
\end{tikzpicture}
\caption{自然堆积视觉筛分需要跨越的三个层次：实例感知、物理尺度与筛分语义。}
\label{fig:three-gaps}
\end{figure}

首先，自然堆积状态下颗粒之间普遍存在接触、局部遮挡以及阴影交叠。此时，视觉系统既需要区分颗粒与背景，又需要恢复相邻颗粒之间并不总是清晰可见的实例边界。若两个颗粒被错误合并为一个实例，其几何尺寸通常会被高估；反之，若单颗粒被过度分割，则会产生偏小的伪颗粒。因此，实例识别误差并不会停留在分割层，而会直接改变后续粒径分布。

其次，实例掩膜本身只提供像素空间中的轮廓。赛题要求输出毫米尺度的长径、短径和粒径分布，并使用 $50\,\mathrm{mm}$ 等物理阈值判断异常，因此必须建立稳定的 pixel--mm 映射。相机工作距离、视角变化、透视畸变和标尺误差都会对这一映射产生影响，而尺度误差又会系统性作用于几乎所有后续长度量。尺度标定因而不是独立的辅助步骤，而是整条物理测量链的共同锚点。

更关键的是，即使每个颗粒的二维轮廓和毫米尺寸均被准确获得，视觉统计结果仍不能直接等同于标准机械筛分结果。一方面，单目图像观察到的是颗粒在当前姿态下的二维投影，而机械筛分反映的是颗粒在运动和姿态变化过程中通过特定筛孔的能力，因此某一个二维长度只能作为筛分粒径的代理量；另一方面，图像天然以“颗粒个数”为统计单位，而机械筛分通过逐级称重得到“颗粒质量”的分布。少量大颗粒在数量统计中占比很低，却可能贡献大部分质量，因此直接按照颗粒数量计算 $D_{10}$、$D_{50}$ 和 $D_{90}$ 会产生系统性的统计口径偏差。

因此，本文将视觉筛分视为一个分层测量与校准问题，而不是单一的实例分割问题：上游需要获得尽可能可靠的颗粒实例，中游需要建立像素到物理几何量的映射，下游则必须进一步将二维视觉几何转换为标准筛分关注的质量级配。为使后续各模块围绕统一目标展开，下一节对这一任务进行形式化定义。

\subsection{问题形式化}
\label{sec:problem-formulation}

给定一幅自然堆积骨料图像 $I$，视觉感知模块首先需要识别其中 $N$ 个可见颗粒，并输出实例掩膜集合
\begin{equation}
\mathcal{M}=\{M_i\}_{i=1}^{N},
\label{eq:instance-set}
\end{equation}
其中 $M_i$ 表示第 $i$ 个颗粒在图像中的像素区域。该定义强调本文关注的是实例级颗粒，而非仅区分“骨料前景”和“背景”的二值语义分割；只有恢复独立实例，后续颗粒级尺寸和形貌计算才具有明确对象。

对每个实例 $M_i$，结合尺度系数将像素轮廓转换为物理几何表征
\begin{equation}
\mathbf{x}_i=(a_i,b_i,A_i,P_i,C_i,S_i,\ldots),
\label{eq:particle-descriptor}
\end{equation}
其中 $a_i$ 和 $b_i$ 分别表示颗粒投影的长轴和短轴，$A_i$ 与 $P_i$ 表示面积和周长，$C_i$ 与 $S_i$ 分别表示圆度和凸度等无量纲形貌指标。在此基础上，系统分别构造筛分等效粒径 $d_i=f(\mathbf{x}_i;\boldsymbol{\theta})$、形貌类别 $c_i=h(\mathbf{x}_i)$ 和异常状态 $z_i=q(\mathbf{x}_i)$。其中，参数 $\boldsymbol{\theta}$ 用于描述二维视觉尺寸与筛分尺寸之间需要由实验校准的系统关系。

赛题最终关注的是场景级粒径级配，而不是彼此独立的单颗粒尺寸。为与机械筛分按质量统计的口径对应，本文进一步为每个颗粒定义质量代理权重
\begin{equation}
w_i=d_i^{\gamma},
\label{eq:intro-mass-weight}
\end{equation}
并基于 $\{d_i,w_i\}_{i=1}^{N}$ 构造累计质量分布 $F(d)$。由此得到的最终粒径输出可以写为
\begin{equation}
\mathcal{Y}=\{D_{10},D_{50},D_{90},P_1,P_2,\ldots,P_K\},
\label{eq:scene-output}
\end{equation}
其中 $P_k$ 表示第 $k$ 个标准筛级对应的预测质量占比。于是，完整任务被分解为三个连续映射：$I\rightarrow\mathcal{M}$ 解决实例感知，$\mathcal{M}\rightarrow\{\mathbf{x}_i\}$ 解决物理测量，$\{\mathbf{x}_i\}\rightarrow\mathcal{Y}$ 解决筛分语义转换。后续方法章节将分别围绕这三个映射展开。

\subsection{本文方案与主要贡献}
\label{sec:contributions}

针对上述问题，本文构建了一套“实例感知--物理测量--筛分语义映射--任务输出”的分层视觉筛分框架。系统首先复用 ImageGrains 2.0\cite{mair2026imagegrains2} 中基于 Cellpose-SAM\cite{pachitariu2025cellposeSAM} 的预训练颗粒分割能力，在尽量避免比赛现场重新标注和训练的条件下恢复自然堆积颗粒实例；随后通过尺度标定和轮廓几何分析获得毫米级颗粒特征；在此基础上，以可解释的筛分等效粒径模型和质量加权模型将二维视觉统计转换为机械筛分可比较的累计质量分布；最后复用同一套颗粒实例和几何表征完成形貌分类、异常检测以及自动化结果输出。

这一设计遵循一个核心原则：对于具有充足预训练数据且主要承担视觉感知功能的环节，优先利用强泛化视觉模型降低现场训练成本；对于直接决定工程筛分口径、且需要通过少量真实筛分批次进行校准的环节，则采用低维、显式和可解释的物理统计模型，而不是将整个任务完全交由端到端黑盒网络学习。这样既利用了基础视觉模型对复杂颗粒边界的表征能力，又保留了从二维测量到标准筛分结果之间可追溯、可校准的关系。

本文的主要工作可以概括为以下三点：
\begin{itemize}[leftmargin=2em]
    \item \textbf{面向自然堆积场景的实例感知与物理测量链。} 基于 ImageGrains 2.0 / Cellpose-SAM 构建颗粒实例分割、尺度恢复和统一几何特征提取流程，为后续粒径、形貌与异常分析提供共享的颗粒级表示。
    \item \textbf{从二维视觉几何到标准筛分质量级配的显式映射。} 通过筛分等效粒径 $d=f(\mathbf{x};\boldsymbol{\theta})$ 与质量代理权重 $w=d^{\gamma}$，区分“视觉数量分布”和“机械筛分质量分布”，并预留基于配对机械筛分真值的参数校准机制。
    \item \textbf{面向竞赛现场约束的一体化自动筛分系统。} 将粒径分析、形貌分类、异常检测、标准化可视化和结果报告集成为可复现软件流程，并以完整现场工作流而非单一模型推理时间评估其部署可行性。
\end{itemize}

上述设计并非由单一模型能力决定，而是由自然堆积条件、标准筛分评价口径和现场部署限制共同推导得到。下一章将从候选技术路线、系统设计原则和总体架构三个层面说明各模块的选型依据及其相互关系。
